\documentclass[12pt,a4paper]{article}

% Packages
\usepackage[utf8]{inputenc}
\usepackage[T1]{fontenc}
\usepackage[french]{babel}
\usepackage{geometry}
\geometry{top=2.5cm, bottom=2.5cm, left=2.8cm, right=2.8cm}

\usepackage{graphicx}
\usepackage{booktabs}
\usepackage{array}
\usepackage{caption}
\usepackage{subcaption}
\usepackage{hyperref}
\usepackage{xcolor}
\usepackage{amsmath}
\usepackage{amssymb}
\usepackage{fancyhdr}
\usepackage{titlesec}
\usepackage{setspace}
\usepackage{listings}
\usepackage{enumitem}
\usepackage{cite}
\usepackage{float}
\usepackage{microtype}
\usepackage{amsopn}

\onehalfspacing

% En-têtes
\pagestyle{fancy}
\fancyhf{}
\rhead{\small TER — Master Informatique 1\textsuperscript{ère} année}
\lhead{\small TOUHAF Aya — LISIC / ULCO}
\cfoot{\thepage}
\renewcommand{\headrulewidth}{0.4pt}

% Formatage sections
\titleformat{\section}{\large\bfseries}{\thesection.}{0.6em}{}[\titlerule]
\titleformat{\subsection}{\normalsize\bfseries}{\thesubsection.}{0.5em}{}
\titleformat{\subsubsection}{\normalsize\itshape}{\thesubsubsection.}{0.5em}{}

% Listings Python
\lstset{
  language=Python,
  basicstyle=\footnotesize\ttfamily,
  keywordstyle=\color{blue!70},
  commentstyle=\color{gray},
  stringstyle=\color{orange!80!black},
  breaklines=true,
  frame=single,
  rulecolor=\color{gray!40},
  backgroundcolor=\color{gray!5},
  numbers=left, numberstyle=\tiny\color{gray},
  xleftmargin=1.5em,
}

\hypersetup{
  colorlinks=true,
  linkcolor=blue!60!black,
  citecolor=green!50!black,
  urlcolor=blue!70,
}

\begin{document}

% PAGE DE TITRE
\begin{titlepage}
\centering
\vspace*{1cm}

{\Large\textsc{Université du Littoral Côte d'Opale}}\\[0.3cm]
{\large Laboratoire d'Informatique Signal et Image de la Côte d'Opale (LISIC)}\\[0.2cm]
{\normalsize Master Informatique — 1\textsuperscript{ère} année}

\vspace{1.5cm}
\rule{\linewidth}{0.8pt}\\[0.4cm]
{\LARGE\bfseries Vers une Intelligence Artificielle Explicable\\[0.2cm]
pour la Surveillance de l'État Physiologique\\[0.2cm]
de l'Athlète de Haut Niveau}\\[0.3cm]
{\large\itshape Espace d'État Fonctionnel, Graphes de Connaissances\\
et Cohérence Sémantique}\\[0.2cm]
\rule{\linewidth}{0.8pt}

\vspace{1.2cm}
{\large Travail d'Étude et de Recherche — 2025--2026}

\vspace{1.5cm}
\begin{tabular}{ll}
\textbf & TOUHAF Aya \\[0.2cm]
                    & \texttt{aya.touhaf@univ-littoral.fr} \\[0.2cm]
\textbf{Encadrant}  & Dr.~Adeel Ahmad \\
                    & \texttt{adeel.ahmad@univ-littoral.fr} \\[0.2cm]
\textbf{Laboratoire} & LISIC — Université du Littoral Côte d'Opale \\[0.2cm]
\textbf{Période}    & Juin -- Juillet 2026 \\
\end{tabular}

\vfill
{\small Rapport de Travail d'Étude et de Recherche — Session 2025--2026}
\end{titlepage}

% RÉSUMÉ
\begin{abstract}
Ce rapport présente les travaux réalisés dans le cadre du TER de Master
Informatique première année, portant sur l'explicabilité des modèles
d'intelligence artificielle appliquée à la surveillance physiologique
d'athlètes de haut niveau équipés de capteurs portables. Face à l'opacité
des modèles d'apprentissage automatique, nous proposons un cadre combinant
un Espace d'État Fonctionnel (FSS), un graphe de connaissances sémantiques
construit automatiquement à partir des valeurs SHAP, et un module de
vérification de cohérence entre les décisions du modèle et le raisonnement
symbolique.

Les expérimentations ont été conduites sur la base SportDB~2.0, en retenant
dix sujets pratiquant la course à pied. Un classifieur Random Forest est
entraîné sur 3\,267 fenêtres temporelles extraites de signaux
cardio-respiratoires : fréquence cardiaque, fréquence respiratoire et
intervalles RR. L'évaluation par leave-one-subject-out révèle une précision
moyenne de 54\,\%, ce qui met en évidence la variabilité inter-individuelle
des données physiologiques. Le score de cohérence IA--graphe de connaissances
atteint en moyenne 0,82, avec 69,4\,\% des prédictions parfaitement validées
par le raisonnement symbolique.

Ces résultats montrent qu'il est possible de relier les prédictions du
modèle à une représentation sémantique cohérente des états physiologiques
observés.

\medskip
\noindent\textbf{Mots-clés :} Explainable AI, XAI, SHAP, Espace d'État
Fonctionnel, Graphe de Connaissances, Random Forest, Wearables, SportDB.
\end{abstract}

\newpage
\tableofcontents
\newpage


\section{Introduction}
\label{sec:intro}

La surveillance en temps réel de l'état physiologique des sportifs de haut
niveau constitue un enjeu important, aussi bien pour la performance que pour
la prévention des blessures. Les dispositifs portables modernes, comme les
montres connectées, les ceintures cardiaques ou les capteurs
cardio-respiratoires, permettent aujourd'hui de collecter en continu des
données biométriques pendant l'entraînement ou la compétition. Bien
analysées, ces données peuvent aider à mieux comprendre l'état d'un athlète,
à repérer des phases de fatigue ou à suivre sa récupération.

L'apprentissage automatique offre des outils intéressants pour exploiter ces
signaux. Cependant, les modèles prédictifs restent souvent difficiles à
interpréter. Dans un contexte sportif, cette limite est importante : un
entraîneur, un préparateur physique ou un médecin du sport ne peut pas se
contenter d'une prédiction brute. Il doit aussi comprendre pourquoi le
système propose une certaine interprétation de l'état de l'athlète.

Cette problématique rejoint le domaine de l'intelligence artificielle
explicable (XAI), dont l'objectif est de rendre les décisions des modèles
plus compréhensibles et plus auditables. Dans des domaines où les décisions
ont un impact humain, comme la santé, le sport ou la sécurité, la capacité à
justifier une recommandation devient presque aussi importante que la
performance du modèle. Cette question est également en lien avec le cadre
réglementaire européen, notamment le Règlement européen sur l'intelligence
artificielle, ou \textit{AI Act}, qui insiste sur la transparence et la
traçabilité des systèmes d'IA à risque élevé~\cite{aiact}.

Les travaux présentés dans ce rapport s'inscrivent dans cette direction.
L'objectif est de construire une première preuve de concept combinant trois
aspects : la modélisation des états physiologiques sous forme d'un espace
d'état fonctionnel, l'extraction d'explications depuis un modèle
d'apprentissage automatique, et la vérification de la cohérence entre ces
explications et une représentation sémantique du domaine.

Ce rapport est organisé comme suit. La section~\ref{sec:etat_art} présente
les principaux éléments liés à l'IA explicable, aux espaces d'état et aux
graphes de connaissances. La section~\ref{sec:methodo} décrit la
méthodologie proposée. La section~\ref{sec:donnees} présente le jeu de
données et le protocole expérimental. Les résultats sont analysés dans la
section~\ref{sec:resultats}, avant une discussion critique et les
perspectives de poursuite du travail.

\section{État de l'art}
\label{sec:etat_art}

\subsection{Intelligence artificielle explicable}

La montée en puissance des modèles d'apprentissage automatique a permis
d'obtenir de bonnes performances dans de nombreux domaines, mais elle a aussi
renforcé le problème de l'opacité des modèles. Linardatos et al.~\cite{linardatos2020}
distinguent deux grandes familles de méthodes d'explicabilité : les méthodes
post-hoc, qui expliquent un modèle déjà entraîné, et les modèles
intrinsèquement interprétables, qui sont lisibles par construction.

Parmi les méthodes post-hoc, LIME~\cite{ribeiro2016} propose d'approximer
localement le comportement d'un modèle complexe par un modèle simple autour
d'une prédiction donnée. Cette approche est intéressante, car elle peut être
utilisée avec différents types de modèles, mais elle peut être sensible aux
paramètres choisis pour générer les perturbations locales.

SHAP~\cite{lundberg2017} repose sur une base théorique plus stable, issue de
la théorie des jeux coopératifs. L'idée est d'attribuer à chaque variable une
contribution dans la prédiction finale. Pour les modèles à base d'arbres,
TreeSHAP~\cite{lundberg2020} permet de calculer ces contributions de façon
efficace. Dans notre travail, SHAP est donc utilisé pour identifier les
variables physiologiques qui influencent le plus les décisions du modèle.

\subsection{Espaces d'état fonctionnels}

La notion d'espace d'état fonctionnel consiste à représenter un système par
un ensemble d'états distincts et par des transitions possibles entre ces
états. Dans notre cas, le système étudié est l'état physiologique de
l'athlète pendant l'effort. Cette représentation est utile, car elle permet
de passer d'une simple prédiction de classe à une lecture plus dynamique du
comportement physiologique.

Pour le protocole RUN, les états suivent naturellement la progression de la
séance : repos, démarrage de l'effort, effort intense, effort soutenu, charge
maximale et récupération. Cette représentation est proche de l'idée des
modèles de transition, où l'état courant du système influence l'état suivant.

\subsection{Graphes de connaissances}

Les graphes de connaissances représentent des concepts et leurs relations
sous forme de nœuds et d'arêtes~\cite{hogan2021}. Ils sont souvent utilisés
pour structurer des connaissances métier et faciliter le raisonnement
symbolique. Dans le domaine médical, Rotmensch et al.~\cite{rotmensch2017}
ont montré qu'il était possible de construire des graphes de connaissances à
partir de données cliniques.

Dans notre cas, le graphe ne sert pas seulement à représenter des concepts
physiologiques. Il sert aussi à relier les variables importantes du modèle
aux états fonctionnels. L'objectif est de vérifier si la décision du modèle
est cohérente avec une lecture sémantique du domaine.


\section{Méthodologie}
\label{sec:methodo}


\subsection{Architecture générale}

L'architecture développée dans ce travail repose sur une chaîne complète
allant des signaux physiologiques jusqu'à l'explication de la décision. Elle
comprend cinq étapes principales :

\begin{enumerate}[leftmargin=1.5em]
  \item extraction de fenêtres temporelles à partir des signaux HR, BR et RR ;
  \item calcul de variables statistiques sur chaque fenêtre ;
  \item entraînement d'un classifieur Random Forest ;
  \item calcul des explications SHAP et extraction de règles ;
  \item construction d'un graphe de connaissances et calcul d'un score de
        cohérence.
\end{enumerate}

L'originalité de cette approche est que le graphe de connaissances n'est pas
construit uniquement à la main. Il est enrichi automatiquement à partir des
explications produites par SHAP. Ainsi, les liens entre variables
physiologiques et états fonctionnels reflètent ce que le modèle apprend dans
les données.

\subsection{Espace d'état fonctionnel}

L'espace d'état fonctionnel est défini à partir du protocole d'exercice
renseigné dans les fichiers \texttt{TrNote.txt}. Pour le sous-ensemble RUN,
six états sont retenus :

\begin{itemize}[leftmargin=1.5em, itemsep=0pt]
  \item $S_0$ — \textit{Resting} : phase de repos avant l'effort ;
  \item $S_1$ — \textit{Effort Initiation} : début de l'effort ;
  \item $S_2$ — \textit{Intense Effort} : effort intense ;
  \item $S_3$ — \textit{Sustained Effort} : effort soutenu ;
  \item $S_4$ — \textit{Peak Load} : charge physiologique maximale ;
  \item $S_5$ — \textit{Recovery} : récupération après l'effort.
\end{itemize}

Les transitions suivent la progression temporelle du protocole :

\[
S_0 \rightarrow S_1 \rightarrow S_2 \rightarrow S_3 \rightarrow S_4 \rightarrow S_5
\]

Cette formalisation permet d'interpréter les erreurs du modèle. Par exemple,
une confusion entre $S_2$ et $S_3$ est moins problématique qu'une confusion
entre $S_0$ et $S_4$, car $S_2$ et $S_3$ sont deux états voisins dans la
dynamique de l'effort.

\subsection{Extraction de features par fenêtre glissante}
\label{subsec:features}

Les signaux HR, BR et RR sont segmentés en fenêtres temporelles. Dans ce
travail, une fenêtre de 20 échantillons est utilisée, avec un pas de 10
échantillons. Cela crée un chevauchement de 50\,\%, ce qui permet de garder
une continuité dans l'analyse du signal.

Pour chaque fenêtre, neuf variables statistiques sont calculées :

\begin{equation}
\mathbf{x} = \bigl[\mu_{HR},\,\sigma_{HR},\,\max_{HR},\;
              \mu_{BR},\,\sigma_{BR},\,\max_{BR},\;
              \mu_{RR},\,\sigma_{RR},\,\min_{RR}\bigr]
\end{equation}

Ces variables résument l'activité cardiaque, respiratoire et la variation
des intervalles RR sur une courte période. Les fenêtres contenant des
valeurs manquantes ou non finies sont supprimées avant l'entraînement.

\subsection{Modèle de classification}

Le modèle utilisé est une Random Forest~\cite{breiman2001}. Ce choix est
adapté à une première preuve de concept, car ce modèle fonctionne bien sur
des données tabulaires, reste robuste avec peu de variables, et permet
d'utiliser facilement TreeSHAP pour l'explicabilité.

Les hyperparamètres utilisés sont les suivants : 200 arbres, une profondeur
maximale de 8, un minimum de 3 échantillons par feuille, et une pondération
des classes pour limiter l'effet du déséquilibre entre états.

\subsection{Construction du graphe de connaissances}

Le graphe de connaissances contient trois types de nœuds :

\begin{itemize}
  \item les états fonctionnels $S_0$ à $S_5$ ;
  \item les indicateurs physiologiques : HR, BR et RR ;
  \item les concepts sémantiques : réponse cardiorespiratoire, charge
        physiologique, risque de fatigue et récupération.
\end{itemize}

Les liens entre variables physiologiques et états sont construits à partir
des valeurs SHAP. Pour chaque classe, on calcule l'importance moyenne des
features, puis on relie l'état au type de signal dominant.

\begin{equation}
\text{shap\_matrix}[c,\, j]
  = \frac{1}{|W|} \sum_{w \in W} \left|\phi_j^{(c)}(w)\right|
\end{equation}

où $\phi_j^{(c)}(w)$ représente la valeur SHAP de la feature $j$ pour la
fenêtre $w$ et la classe $c$.

\subsection{Score de cohérence IA--KG}

Pour chaque fenêtre, on compare la prédiction du modèle avec l'inférence
obtenue par les règles extraites automatiquement des arbres. Le score de
cohérence est défini ainsi :

\begin{equation}
\text{CS}(\hat{y}_{\text{IA}},\,\hat{y}_{\text{KG}}) =
\begin{cases}
1{,}0 & \text{si } \hat{y}_{\text{IA}} = \hat{y}_{\text{KG}} \\
0{,}7 & \text{si les deux états sont voisins} \\
0{,}4 & \text{si les deux états sont séparés par deux transitions} \\
0{,}0 & \text{sinon}
\end{cases}
\end{equation}

Ce choix tient compte de la structure du FSS. Une erreur entre deux états
voisins n'a pas le même sens qu'une incohérence complète entre deux états
éloignés.


\section{Données et protocole expérimental}
\label{sec:donnees}


\subsection{Dataset SportDB 2.0}

Le jeu de données utilisé est SportDB~2.0~\cite{sportdb2}. Il contient des
enregistrements physiologiques collectés sur 130 sujets pendant la pratique
de 11 sports. Les signaux ont été acquis avec différents capteurs portables,
notamment BioHarness, Polar et KardiaMobile.

Dans ce travail, nous avons choisi le sous-ensemble \textbf{RUN}, car il
présente un protocole structuré et comparable entre les sujets. Les dix
sujets disponibles, de S1 à S10, ont été utilisés.

\begin{table}[H]
\centering
\caption{Caractéristiques du dataset après prétraitement.}
\label{tab:dataset}
\begin{tabular}{lc}
\toprule
\textbf{Paramètre} & \textbf{Valeur} \\
\midrule
Nombre de sujets             & 10 \\
Nombre de fenêtres totales   & 3\,267 \\
Taille de fenêtre            & 20 échantillons \\
Pas de glissement            & 10 échantillons \\
Nombre de features           & 9 \\
Nombre de classes            & 6 \\
\bottomrule
\end{tabular}
\end{table}

\subsection{Répartition des fenêtres par état}

La répartition des fenêtres n'est pas parfaitement équilibrée, car les
phases du protocole n'ont pas toutes la même durée.

\begin{table}[H]
\centering
\caption{Nombre de fenêtres par état FSS.}
\label{tab:windows}
\begin{tabular}{lcc}
\toprule
\textbf{État} & \textbf{Fenêtres} & \textbf{\%} \\
\midrule
$S_0$ — Resting           & 617  & 18,9\,\% \\
$S_1$ — Effort Initiation & 334  & 10,2\,\% \\
$S_2$ — Intense Effort    & 554  & 17,0\,\% \\
$S_3$ — Sustained Effort  & 390  & 11,9\,\% \\
$S_4$ — Peak Load         & 754  & 23,1\,\% \\
$S_5$ — Recovery          & 618  & 18,9\,\% \\
\midrule
\textbf{Total}            & \textbf{3\,267} & \textbf{100\,\%} \\
\bottomrule
\end{tabular}
\end{table}

\subsection{Protocole de validation}

Deux protocoles de validation ont été utilisés. La validation croisée
stratifiée en 5 plis donne une estimation de la performance lorsque les
données sont mélangées entre sujets. Le protocole leave-one-subject-out
(LOSO) est plus exigeant : à chaque itération, toutes les fenêtres d'un
sujet sont retirées de l'entraînement et utilisées pour le test. Ce type de
validation est particulièrement adapté aux données physiologiques et à la
reconnaissance d'activités humaines, car il permet d'évaluer la capacité
d'un modèle à généraliser à des individus jamais observés pendant
l'entraînement~\cite{loso_cv}.


\section{Résultats}
\label{sec:resultats}


\subsection{Performance du classifieur}

\subsubsection{Matrice de confusion}

La figure~\ref{fig:confusion} présente la matrice de confusion obtenue. Les
états $S_0$ et $S_5$ sont bien reconnus, car ils correspondent à des profils
physiologiques assez distincts. Les principales confusions apparaissent
entre $S_2$, $S_3$ et $S_4$, ce qui est cohérent avec la proximité de ces
phases dans le déroulement de l'effort.

\begin{figure}[H]
  \centering
  \includegraphics[width=0.78\linewidth]{SportDB2/RUN/outputs/1_confusion_matrix}
  \caption{Matrice de confusion du classifieur Random Forest.}
  \label{fig:confusion}
\end{figure}

\subsubsection{Validation croisée et généralisation}

\begin{table}[H]
\centering
\caption{Résultats de validation du classifieur.}
\label{tab:scores}
\begin{tabular}{lcc}
\toprule
\textbf{Protocole} & \textbf{Accuracy moyenne} & \textbf{Écart-type} \\
\midrule
Cross-validation 5-fold & 0{,}717 & \pm 0{,}022 \\
Leave-One-Subject-Out   & 0{,}540 & \pm 0{,}103 \\
\bottomrule
\end{tabular}
\end{table}

L'écart entre la validation croisée et le LOSO montre que le modèle apprend
correctement les états dans l'ensemble des données, mais que la
généralisation à un sujet jamais vu reste plus difficile. Ce résultat est
attendu dans des données physiologiques, car deux athlètes peuvent présenter
des réponses cardio-respiratoires différentes pour un effort similaire.

\subsection{Analyse SHAP}

\begin{figure}[H]
  \centering
  \includegraphics[width=0.85\linewidth]{SportDB2/RUN/outputs/2_shap_summary}
  \caption{Résumé des valeurs SHAP par classe.}
  \label{fig:shap_summary}
\end{figure}

La figure~\ref{fig:shap_summary} montre que les variables liées à la
fréquence cardiaque, en particulier \texttt{max\_HR} et \texttt{mean\_HR},
sont les plus importantes. Les variables liées aux intervalles RR jouent
également un rôle important. Ce résultat est cohérent : quand l'intensité
de l'effort augmente, la fréquence cardiaque augmente et les intervalles RR
diminuent.

\begin{figure}[H]
  \centering
  \includegraphics[width=\linewidth]{SportDB2/RUN/outputs/4_feature_importance}
  \caption{Comparaison des importances MDI et SHAP.}
  \label{fig:feature_importance}
\end{figure}

La comparaison entre MDI et SHAP confirme que les variables cardiaques
dominent la décision du modèle. Le fait que deux méthodes différentes
mettent en avant les mêmes variables renforce la confiance dans
l'interprétation.

\subsection{Graphe de connaissances}

\begin{figure}[H]
  \centering
  \includegraphics[width=\linewidth]{SportDB2/RUN/outputs/3_knowledge_graph}
  \caption{Graphe de connaissances construit à partir des valeurs SHAP.}
  \label{fig:kg}
\end{figure}

Le graphe obtenu relie les signaux physiologiques aux états fonctionnels et
aux concepts sémantiques. Les liens verts représentent les correspondances
dérivées de SHAP. On retrouve notamment le rôle central de la fréquence
cardiaque dans les états d'effort, ainsi que le rôle des intervalles RR dans
la distinction entre repos et effort.

\subsection{Cohérence IA--Knowledge Graph}

\begin{figure}[H]
  \centering
  \includegraphics[width=\linewidth]{SportDB2/RUN/outputs/5_loso_and_consistency}
  \caption{Accuracy LOSO par sujet et cohérence IA--KG par état.}
  \label{fig:loso_consistency}
\end{figure}

\begin{figure}[H]
  \centering
  \includegraphics[width=\linewidth]{SportDB2/RUN/outputs/6_consistency_and_top_features}
  \caption{Distribution du score de cohérence et feature dominante par état.}
  \label{fig:consistency_distribution}
\end{figure}

Le score de cohérence moyen est de \textbf{0,82}. Sur les 3\,267 fenêtres,
69,4\,\% sont parfaitement cohérentes avec le raisonnement symbolique,
24,7\,\% sont partiellement cohérentes, et 5,9\,\% sont incohérentes. Ce
résultat montre que la majorité des décisions du modèle peuvent être reliées
à une explication compatible avec le graphe de connaissances.

\begin{table}[H]
\centering
\caption{Score de cohérence IA--KG par état FSS.}
\label{tab:consistency_state}
\begin{tabular}{lcc}
\toprule
\textbf{État} & \textbf{Score moyen} & \textbf{Interprétation} \\
\midrule
$S_0$ — Resting           & 0,851 & Bon \\
$S_1$ — Effort Initiation & 0,674 & Modéré \\
$S_2$ — Intense Effort    & 0,842 & Bon \\
$S_3$ — Sustained Effort  & 0,834 & Bon \\
$S_4$ — Peak Load         & 0,736 & Acceptable \\
$S_5$ — Recovery          & 0,936 & Excellent \\
\midrule
\textbf{Moyenne}          & \textbf{0,819} & Bon \\
\bottomrule
\end{tabular}
\end{table}

\subsection{Exemple d'explication générée}

Le système produit également une explication textuelle pour chaque fenêtre.
Un exemple typique est le suivant :

\begin{quote}
\textit{Le modèle prédit l'état $S_0$ avec une confiance élevée. La variable
la plus contributive est \texttt{mean\_HR}. L'inférence par règles conclut
également à $S_0$. Le score de cohérence est donc égal à 1,0.}
\end{quote}

Ce type d'explication permet de relier une décision numérique à une lecture
plus compréhensible de l'état physiologique.


\section{Discussion}
\label{sec:discussion}


\subsection{Apports du travail}

Ce travail a permis de construire un premier prototype complet répondant aux
trois axes du sujet : formalisation d'un espace d'état fonctionnel,
extraction d'explications et de règles, puis évaluation de la cohérence avec
un graphe de connaissances.

Le point le plus intéressant est que le graphe est enrichi à partir des
explications SHAP. Cela évite de construire toutes les relations à la main
et permet de garder un lien direct entre ce que le modèle apprend et la
représentation sémantique utilisée pour l'explication.

\subsection{Limites}

La principale limite est que cette première version utilise seulement le
sous-ensemble RUN de SportDB~2.0. Le nombre de sujets reste donc limité, et
le modèle n'est pas encore généralisé à d'autres sports. De plus, le Random
Forest travaille sur des features statistiques extraites de fenêtres, mais
ne modélise pas directement la dynamique temporelle complète des signaux.

Une autre limite concerne les règles extraites. Elles donnent une
approximation lisible du comportement du modèle, mais elles ne doivent pas
être interprétées comme des règles médicales ou physiologiques validées par
des experts.

\subsection{Lien avec l'explicabilité réglementaire}

Le score de cohérence proposé peut être vu comme une première manière de
documenter la relation entre décision automatique et raisonnement
symbolique. Cette idée rejoint les exigences de transparence et de
traçabilité mises en avant dans l'AI Act~\cite{aiact}. Même si le système
développé ici reste un prototype de recherche, cette direction semble
pertinente pour des applications futures dans des contextes sensibles.


\section{Conclusion et perspectives}
\label{sec:conclusion}


Ce TER a permis de développer une première architecture d'IA explicable pour
le suivi physiologique d'athlètes. Le système combine un espace d'état
fonctionnel, un modèle Random Forest, des explications SHAP, des règles
extraites automatiquement, un graphe de connaissances et un score de
cohérence IA--KG.

Les résultats obtenus sur le sous-ensemble RUN de SportDB~2.0 sont
encourageants. La performance LOSO moyenne de 54\,\% montre que la
généralisation inter-sujets reste difficile, mais le score de cohérence
moyen de 0,82 montre que les décisions du modèle peuvent souvent être
expliquées de manière compatible avec le graphe sémantique.

La prochaine étape consistera à étendre cette approche au dataset SportDB~2.0
complet, afin de tester la méthode sur plusieurs sports. Une autre
perspective importante est le passage à des modèles de deep learning
séquentiels, comme les CNN, LSTM ou modèles avec attention, capables
d'exploiter directement la dynamique temporelle des signaux.

\bigskip
\noindent\textbf{Remerciements.} Je remercie le Dr.~Adeel Ahmad pour son
encadrement, sa disponibilité et ses retours tout au long de ce projet.

\bibliographystyle{plain}
\begin{thebibliography}{99}

\bibitem{linardatos2020}
P.~Linardatos, V.~Papastefanopoulos, and S.~Kotsiantis.
\newblock Explainable AI: A review of machine learning interpretability methods.
\newblock \emph{Entropy}, 23(1):18, 2021.
\newblock \url{https://www.mdpi.com/1099-4300/23/1/18}

\bibitem{ribeiro2016}
M.~T.~Ribeiro, S.~Singh, and C.~Guestrin.
\newblock ``Why Should I Trust You?'': Explaining the Predictions of Any Classifier.
\newblock KDD, 2016.
\newblock \url{https://arxiv.org/abs/1602.04938}

\bibitem{lundberg2017}
S.~M.~Lundberg and S.-I.~Lee.
\newblock A Unified Approach to Interpreting Model Predictions.
\newblock NeurIPS, 2017.
\newblock \url{https://arxiv.org/abs/1705.07874}

\bibitem{lundberg2020}
S.~M.~Lundberg et al.
\newblock Explainable AI for Trees: From Local Explanations to Global Understanding.
\newblock 2019.
\newblock \url{https://arxiv.org/abs/1905.04610}

\bibitem{breiman2001}
L.~Breiman.
\newblock Random Forests.
\newblock \emph{Machine Learning}, 45:5--32, 2001.
\newblock \url{https://www.stat.berkeley.edu/~breiman/randomforest2001.pdf}

\bibitem{hogan2021}
A.~Hogan et al.
\newblock Knowledge Graphs.
\newblock \emph{ACM Computing Surveys}, 54(4), 2021.
\newblock \url{https://arxiv.org/abs/2003.02320}

\bibitem{rotmensch2017}
M.~Rotmensch, Y.~Halpern, A.~Tlimat, S.~Horng, and D.~Sontag.
\newblock Learning a Health Knowledge Graph from Electronic Medical Records.
\newblock \emph{Scientific Reports}, 7:5994, 2017.
\newblock \url{https://www.nature.com/articles/s41598-017-05778-z}

\bibitem{claudino2019}
J.~G.~Claudino et al.
\newblock Current Approaches to the Use of Artificial Intelligence for Injury Risk Assessment and Performance Prediction in Team Sports.
\newblock \emph{Sports Medicine -- Open}, 5:28, 2019.
\newblock \url{https://link.springer.com/article/10.1186/s40798-019-0202-3}

\bibitem{sportdb2}
L.~Burattini, S.~Romagnoli, A.~Sbrollini, M.~Morettini, A.~Nocera,
D.~Bondi, T.~Pietrangelo, and V.~Verratti.
\newblock Sport DB 2.0.
\newblock Mendeley Data, Version 2, 2024.
\newblock \url{https://data.mendeley.com/datasets/kzkjkt7mx2}

\bibitem{sportdb2_cinc}
S.~Romagnoli, A.~Sbrollini, A.~Nocera, M.~Morettini, D.~Bondi,
T.~Pietrangelo, V.~Verratti, and L.~Burattini.
\newblock Sport DB 2.0: A New Database of Data Acquired by Wearable and Portable Devices While Practicing Sport.
\newblock Computing in Cardiology, 2023.
\newblock \url{https://cinc.org/archives/2023/pdf/CinC2023-067.pdf}

\bibitem{loso_cv}
A.~Dehghani, T.~Glatard, and E.~Shihab.
\newblock Subject Cross Validation in Human Activity Recognition.
\newblock arXiv, 2019.
\newblock \url{https://arxiv.org/abs/1904.02666}

\bibitem{aiact}
European Parliament and Council of the European Union.
\newblock Regulation (EU) 2024/1689 laying down harmonised rules on artificial intelligence.
\newblock Official Journal of the European Union, 2024.
\newblock \url{https://eur-lex.europa.eu/eli/reg/2024/1689/oj/eng}

\end{thebibliography}

\end{document}